%% wg21org-preamble.tex --- preamble for WG21 papers exported by ox-wg21latex.el
%%
%% The exporter copies this into every paper, after the paper's own
%% LATEX_HEADER lines, so an exported paper compiles with nothing but
%% TeX Live.  Anything the paper already defined, such as the wording
%% macros from common.tex, is kept: every definition here is provided,
%% not imposed.

\makeatletter

\usepackage{xcolor}
\usepackage[normalem]{ulem}
\usepackage{array}
\usepackage{booktabs}
\usepackage{longtable}
\usepackage{hyperref}
\usepackage{iftex}

%% Nothing should run into the margin: URLs break anywhere, identifiers
%% after an underscore, and a paragraph that cannot be set well is set
%% loosely instead of overfull.
\usepackage{xurl}
\usepackage{underscore}
\setlength{\emergencystretch}{3em}

%% pdfLaTeX knows only the Unicode characters declared to it.  A
%% zero-width space is a place a line may break.
\ifPDFTeX
  \DeclareUnicodeCharacter{200B}{\hspace{0pt}}
\fi
%% Links as coloured text, not boxes: cross references and the
%% contents in the text colour, links out in blue.
\hypersetup{colorlinks=true,
             linkcolor=.,
             citecolor=.,
             urlcolor=[rgb]{0.04,0.34,0.82}}

%% Wording changes: [[insert:][...]], [[delete:][...]], _underline_,
%% +strike-through+, and the addedblock and removedblock blocks.
\providecolor{addclr}{rgb}{0,.45,.15}
\providecolor{remclr}{rgb}{.75,0,0}
\providecommand{\added}[1]{\textcolor{addclr}{\uline{#1}}}
\providecommand{\removed}[1]{\textcolor{remclr}{\sout{#1}}}
\@ifundefined{addedblock}{\newenvironment{addedblock}{\color{addclr}}{}}{}
\@ifundefined{removedblock}{\newenvironment{removedblock}{\color{remclr}}{}}{}

%% Title block: the exporter writes the title and a table of the
%% document's metadata, rather than using \maketitle, which common.tex
%% and document classes each define their own way.
\providecommand{\wgmetalabel}[1]{\textsf{\small #1}}
\providecommand{\wgmetavalue}[1]{\textsf{\small #1}}

%% A special block, #+begin_NAME, is the environment NAME when the
%% document has one; otherwise its contents are set as they are, with a
%% warning, rather than stopping the build.
\newenvironment{wgblock}[1]
  {\@ifundefined{#1}
    {\PackageWarning{wg21org}{No environment `#1'; its block is set as plain text}%
     \def\wg@blockend{}}
    {\def\wg@blockend{\end{#1}}\begin{#1}}}
  {\wg@blockend}

%% Footnotes.  The working draft's macros (stdtex/macros.tex, from
%% common.tex) replace the \footnote command with a footnote
%% environment; use whichever the document has.
\providecommand{\wgfootnote}[1]{%
  \@ifundefined{endfootnote}{\footnote{#1}}{\begin{footnote}#1\end{footnote}}}

%% The table of contents, without memoir's own entry for it in itself;
%% only memoir has the starred form.
\providecommand{\wgtableofcontents}{%
  \@ifclassloaded{memoir}{\tableofcontents*}{\tableofcontents}}

%% A section's link to its line in the Org source, in the margin.
%% It takes no width, since margins can be narrower than the word.
\providecommand{\wgsourcelink}[1]{%
  \marginpar{\makebox[0pt][l]{\textsf{\scriptsize\href{#1}{source}}}}}

%% Comparison tables: before and after columns of equal width, as a
%% longtable so a long table breaks between rows, with its header row
%% repeated on each page.
\newlength{\wgcmptblgap}
\setlength{\wgcmptblgap}{1.5em}
\newlength{\wgcmptblcol}
\newenvironment{wgcmptbl}
  {\setlength{\wgcmptblcol}{\dimexpr(\linewidth-\wgcmptblgap)/2\relax}%
   \setlength{\LTpre}{\medskipamount}%
   \setlength{\LTpost}{\medskipamount}%
   \let\wg@pagebreak\pagebreak
   \let\wg@nopagebreak\nopagebreak
   \let\wg@newpage\newpage
   \begin{longtable}{@{}p{\wgcmptblcol}@{\hspace{\wgcmptblgap}}p{\wgcmptblcol}@{}}}
  {\end{longtable}}
%% A cell.  longtable redefines \pagebreak, \nopagebreak and \newpage
%% to act between rows, and a breakable code box calls them from inside
%% the cell; give the cell back the ordinary meanings, and a minipage
%% to keep it in one piece.
\newenvironment{wgcmptblcell}
  {\let\pagebreak\wg@pagebreak
   \let\nopagebreak\wg@nopagebreak
   \let\newpage\wg@newpage
   \begin{minipage}[t]{\linewidth}}
  {\end{minipage}}

\makeatother
